\subsection{Union-Find / DSU con Path Halving (componentes conexas)}
\label{alg:4-9}

\graybold{Preguntas guía:}

\begin{itemize}
\item ¿Me piden cuántos grupos/islas/componentes separados hay en un grafo no dirigido, o si dos elementos están en el mismo grupo?
\item ¿Las aristas llegan de forma incremental (una por una) y solo necesito consultar conectividad, sin nunca eliminar aristas?
\item ¿Puedo tolerar aristas repetidas y lazos sin que el conteo se rompa?
\item ¿Prefiero evitar construir listas de adyacencia y hacer DFS/BFS explícito?
\end{itemize}

\graybold{Utilidad:} Mantiene una partición de elementos en conjuntos disjuntos soportando dos operaciones: unir los conjuntos de dos elementos, y averiguar a qué conjunto pertenece un elemento. Contar componentes conexas es su aplicación más directa, pero la estructura es la base de Kruskal (árbol de recubrimiento mínimo), detección de ciclos en grafos no dirigidos, agrupamiento por equivalencias y problemas de conectividad incremental en general. A diferencia de DFS/BFS, no requiere materializar el grafo: procesa las aristas al vuelo, una por una.

\graybold{Complejidad:} \bigO{M \log N} en el peor caso para procesar $M$ aristas sobre $N$ nodos, porque aquí se usa solo compresión de caminos. Si además se uniera por tamaño o rango bajaría a \bigO{M \alpha(N)}, donde $\alpha$ es la inversa de la función de Ackermann (menor que 5 para cualquier entrada imaginable), es decir prácticamente \bigO{M}.

\graybold{Fundamento teórico:} Cada conjunto se representa como un árbol donde cada nodo apunta a su padre y la RAÍZ (el nodo que es su propio padre) actúa como identificador canónico del conjunto. Dos elementos pertenecen al mismo conjunto exactamente cuando su búsqueda de raíz termina en el mismo nodo, así que el número de componentes conexas es el número de nodos que son su propia raíz. Unir dos conjuntos cuesta una sola asignación: se hace que la raíz de uno apunte a la raíz del otro. La eficiencia depende de mantener los árboles poco profundos, y para eso se usa COMPRESIÓN DE CAMINOS: durante la búsqueda de la raíz, cada nodo visitado se reengancha más arriba en el árbol, de modo que las consultas futuras sobre ese nodo (y sobre todos sus descendientes en el camino) sean más cortas. La variante implementada aquí, \emph{path halving}, hace que cada nodo apunte a su abuelo en lugar de a su padre mientras sube, lo que reduce la profundidad a la mitad en cada pasada con un solo bucle y sin recursión ni segunda pasada. Con solo esta compresión, el costo amortizado por operación es a lo sumo logarítmico; combinada con unión por tamaño o rango quedaría acotado por $\alpha(N)$, una función que crece tan lentamente que es constante en la práctica. Nótese que la estructura es naturalmente inmune a aristas repetidas y lazos: si $a$ y $b$ ya comparten raíz, la unión simplemente no hace nada, así que la entrada puede contener \texttt{CE} y \texttt{EC} (como en el ejemplo) o incluso \texttt{AA} sin afectar el conteo.

\graybold{Ejercicio aplicado}

(UVA 459 — Graph Connectivity) Dado el nombre del nodo mayor del grafo (una letra mayúscula, que implica los nodos desde \texttt{A} hasta esa letra) y una lista de aristas no dirigidas expresadas como pares de letras, contar cuántos subgrafos conexos maximales (componentes conexas) tiene el grafo.

\graybold{I/O:} La entrada separa los casos con LÍNEAS EN BLANCO, lo que normalmente forzaría a leer línea por línea; sin embargo, aquí la estructura por líneas es redundante: el nodo mayor siempre es un token de 1 carácter y cada arista es un token de 2 caracteres. Eso permite usar lectura total + tokenización (patrón 2) y distinguir el inicio de cada caso por la LONGITUD del token, consumiendo aristas mientras midan 2 — el parseo se vuelve inmune a cualquier cantidad de líneas en blanco, espacios extra o retornos de carro. Se opera directamente sobre \texttt{bytes}, restando 64 al código ASCII para mapear \texttt{'A'}$\to$1 hasta \texttt{'Z'}$\to$26 (a lo sumo 26 nodos, así que el rendimiento es irrelevante). La salida exige una línea en blanco entre casos consecutivos, lo que se logra con \texttt{"\textbackslash n\textbackslash n".join(out)}. Verificado contra el caso de ejemplo y contrastado con una fuerza bruta de BFS sobre listas de adyacencia en 120 tandas aleatorias (incluyendo lazos \texttt{AA} y aristas duplicadas), más casos borde de un solo nodo y de casos sin ninguna arista, sin discrepancias.

\graybold{Código solución}

\begin{lstlisting}[language=Python]
import sys

def solve():
    data = sys.stdin.buffer.read().split()
    idx = 0
    t = int(data[idx]); idx += 1
    out = []
    for _ in range(t):
        # el nodo mayor es un token de 1 caracter; las aristas son de 2
        n = data[idx][0] - 64          # 'A'=65 -> 1 nodo ... 'Z'=90 -> 26 nodos
        idx += 1
        padre = list(range(n + 1))

        def find(x):
            while padre[x] != x:
                padre[x] = padre[padre[x]]   # path halving: apunta al abuelo
                x = padre[x]
            return x

        # consume aristas mientras el token siguiente tenga 2 caracteres:
        # asi el parseo ignora por completo las lineas en blanco del enunciado
        while idx < len(data) and len(data[idx]) == 2:
            tok = data[idx]; idx += 1
            a = find(tok[0] - 64)
            b = find(tok[1] - 64)
            if a != b:                  # si ya comparten raiz (arista repetida
                padre[a] = b            # o lazo) no se hace nada
        # cada nodo que es su propia raiz representa una componente
        out.append(str(sum(1 for v in range(1, n + 1) if find(v) == v)))
    sys.stdout.write("\n\n".join(out) + "\n")

solve()
\end{lstlisting}
